\clearpage
\begin{trapbox}
\textbf{\textcolor{CTrap}{误区一：换开口方向，还照抄 $x_0+p/2$。}}
左右开口量水平距离，上下开口量竖直距离；向左、向下时相应坐标非正，须取相反数。
\begin{center}\begin{tikzpicture}[>=stealth,font=\small]
\draw[rounded corners=3pt,gray!45] (0,0) rectangle (14.2,13.6);
\node at (7.1,13.05) {四种开口，同一件事：焦半径等于到准线的垂直距离};
\node at (3.6,12.35) {向右：$y^2=2px$};
\begin{scope}[shift={(3.6,9.8)},scale=.64]
\draw[->,gray] (-2.5,0)--(2.5,0) node[right] {$x$};
\draw[->,gray] (0,-2.5)--(0,2.5) node[above] {$y$};
\begin{scope}[rotate=0]
\draw[teal!75!black,thick,domain=-2.45:2.45,samples=90] plot ({\x*\x/4},\x);
\draw[red!75!black,thick] (-1,-2.5)--(-1,2.5);
\draw[orange!85!black,very thick] (-1,2)--(1,2);
\draw[blue!75!black,very thick] (1,0)--(1,2);
\draw[red!75!black] (-1,1.7)--(-.7,1.7)--(-.7,2);
\fill (1,2) circle(2pt);
\fill (-1,2) circle(2pt);
\fill (1,0) circle(2pt);
\fill (0,0) circle(2pt);
\end{scope}
\node[below right=3pt] at (1.0,2.0) {$P$};
\node[above left=3pt] at (-1.0,2.0) {$Q$};
\node[below right=3pt] at (1.0,0.0) {$F$};
\node[below left=3pt] at (0.0,0.0) {$O$};
\end{scope}
\node at (3.6,7.5) {$\displaystyle PF=x_0+\frac p2$};
\node at (10.6,12.35) {向左：$y^2=-2px$};
\begin{scope}[shift={(10.6,9.8)},scale=.64]
\draw[->,gray] (-2.5,0)--(2.5,0) node[right] {$x$};
\draw[->,gray] (0,-2.5)--(0,2.5) node[above] {$y$};
\begin{scope}[rotate=180]
\draw[teal!75!black,thick,domain=-2.45:2.45,samples=90] plot ({\x*\x/4},\x);
\draw[red!75!black,thick] (-1,-2.5)--(-1,2.5);
\draw[orange!85!black,very thick] (-1,2)--(1,2);
\draw[blue!75!black,very thick] (1,0)--(1,2);
\draw[red!75!black] (-1,1.7)--(-.7,1.7)--(-.7,2);
\fill (1,2) circle(2pt);
\fill (-1,2) circle(2pt);
\fill (1,0) circle(2pt);
\fill (0,0) circle(2pt);
\end{scope}
\node[above left=3pt] at (-1.0,-2.0) {$P$};
\node[below right=3pt] at (1.0,-2.0) {$Q$};
\node[above left=3pt] at (-1.0,0.0) {$F$};
\node[below right=3pt] at (-0.0,0.0) {$O$};
\end{scope}
\node at (10.6,7.5) {$\displaystyle PF=-x_0+\frac p2$};
\node at (3.6,6.35) {向上：$x^2=2py$};
\begin{scope}[shift={(3.6,3.8)},scale=.64]
\draw[->,gray] (-2.5,0)--(2.5,0) node[right] {$x$};
\draw[->,gray] (0,-2.5)--(0,2.5) node[above] {$y$};
\begin{scope}[rotate=90]
\draw[teal!75!black,thick,domain=-2.45:2.45,samples=90] plot ({\x*\x/4},\x);
\draw[red!75!black,thick] (-1,-2.5)--(-1,2.5);
\draw[orange!85!black,very thick] (-1,2)--(1,2);
\draw[blue!75!black,very thick] (1,0)--(1,2);
\draw[red!75!black] (-1,1.7)--(-.7,1.7)--(-.7,2);
\fill (1,2) circle(2pt);
\fill (-1,2) circle(2pt);
\fill (1,0) circle(2pt);
\fill (0,0) circle(2pt);
\end{scope}
\node[below left=3pt] at (-2.0,1.0) {$P$};
\node[below left=3pt] at (-2.0,-1.0) {$Q$};
\node[above right=3pt] at (0.0,1.0) {$F$};
\node[below right=3pt] at (0.0,0.0) {$O$};
\end{scope}
\node at (3.6,1.5) {$\displaystyle PF=y_0+\frac p2$};
\node at (10.6,6.35) {向下：$x^2=-2py$};
\begin{scope}[shift={(10.6,3.8)},scale=.64]
\draw[->,gray] (-2.5,0)--(2.5,0) node[right] {$x$};
\draw[->,gray] (0,-2.5)--(0,2.5) node[above] {$y$};
\begin{scope}[rotate=270]
\draw[teal!75!black,thick,domain=-2.45:2.45,samples=90] plot ({\x*\x/4},\x);
\draw[red!75!black,thick] (-1,-2.5)--(-1,2.5);
\draw[orange!85!black,very thick] (-1,2)--(1,2);
\draw[blue!75!black,very thick] (1,0)--(1,2);
\draw[red!75!black] (-1,1.7)--(-.7,1.7)--(-.7,2);
\fill (1,2) circle(2pt);
\fill (-1,2) circle(2pt);
\fill (1,0) circle(2pt);
\fill (0,0) circle(2pt);
\end{scope}
\node[above right=3pt] at (2.0,-1.0) {$P$};
\node[above right=3pt] at (2.0,1.0) {$Q$};
\node[below left=3pt] at (-0.0,-1.0) {$F$};
\node[above left=3pt] at (0.0,-0.0) {$O$};
\end{scope}
\node at (10.6,1.5) {$\displaystyle PF=-y_0+\frac p2$};
\node at (7.1,.65) {四图均取 $p=2$；橙段 $PQ$ 与蓝段 $PF$ 都长 $2$，红线为准线。};
\end{tikzpicture}
\end{center}
四式的共同结构都是 $PF=t+p/2$，$t\ge0$。
不能把有正负号的坐标直接当作长度；主图的分段也说明向右应当加 $p/2$，不能减。
\end{trapbox}
\clearpage
\begin{trapbox}
\textbf{\textcolor{CTrap}{误区二：任意点都能用焦半径公式。}}
两图用同一条曲线、同一焦点、同一尺度，且两点横坐标都是 $3$。
\begin{center}\begin{tikzpicture}[>=stealth,font=\small]
\draw[rounded corners=3pt,gray!45] (0,0) rectangle (14.2,9.6);
\node at (7.1,9.05) {横坐标相同，也必须先检查点是否在抛物线上};
\begin{scope}[shift={(2.2,3.6)},scale=.78]
\draw[->,gray] (-1.6,0)--(4.6,0) node[right] {$x$};
\draw[->,gray] (0,-1.2)--(0,4.9) node[above] {$y$};
\draw[red!75!black,thick] (-1,-1.1)--(-1,4.8);
\draw[teal!75!black,thick,domain=-1:4.1,samples=90] plot ({\x*\x/4},\x);
\fill (0,0) circle(1.7pt) node[below left=2pt] {$O$};
\fill (1,0) circle(2pt) node[below=3pt] {$F$};
\draw[orange!85!black,very thick] (-1,{sqrt(12)})--(3,{sqrt(12)});
\draw[blue!75!black,very thick] (1,0)--(3,{sqrt(12)});
\fill (3,{sqrt(12)}) circle(2pt) node[above right=2pt] {$P$};
\fill (-1,{sqrt(12)}) circle(2pt) node[left=2pt] {$Q$};
\draw[gray,dashed] (3,0)--(3,{sqrt(12)});
\node[below] at (3,0) {$3$};
\end{scope}
\node[align=center] at (3.2,1.55) {$P(3,2\sqrt3)$ 在 $y^2=4x$ 上\\[6pt]$PF=PQ=3+1=4$};
\begin{scope}[shift={(9.2,3.6)},scale=.78]
\draw[->,gray] (-1.6,0)--(4.6,0) node[right] {$x$};
\draw[->,gray] (0,-1.2)--(0,4.9) node[above] {$y$};
\draw[red!75!black,thick] (-1,-1.1)--(-1,4.8);
\draw[teal!75!black,thick,domain=-1:4.1,samples=90] plot ({\x*\x/4},\x);
\fill (0,0) circle(1.7pt) node[below left=2pt] {$O$};
\fill (1,0) circle(2pt) node[below=3pt] {$F$};
\draw[orange!85!black,very thick] (-1,0)--(3,0);
\draw[blue!75!black,very thick] (1,0)--(3,0);
\draw[<->,orange!85!black] (-1,-1.6)--(3,-1.6) node[midway,below] {$RQ=4$};
\draw[<->,blue!75!black] (1,1)--(3,1) node[midway,above] {$RF=2$};
\fill (3,0) circle(2pt) node[below right=3pt] {$R(3,0)$};
\fill (-1,0) circle(2pt) node[below left=2pt] {$Q$};
\end{scope}
\node[align=center] at (10.2,1.05) {$R(3,0)$ 不在 $y^2=4x$ 上\\[6pt]$RF\ne RQ$，不能套公式};
\end{tikzpicture}
\end{center}
\textbf{为什么错：}$R$ 不满足 $y^2=4x$，所以没有定义所保证的 $RF=RQ$。
\textbf{正确替代：}先代入检验点在曲线上；不在曲线上，就使用两点距离公式。

\medskip
\textbf{\textcolor{CTrap}{误区三：把 $p/2$ 当作半通径，或读错方程系数。}}
由参数图，半通径是 $p$；公式常数为 $p/2$。
若 $y^2=kx$（$k>0$），则 $k=2p=4a$，公式常数为 $k/4=a$。
通径端点 $x_0=p/2$ 的焦半径是 $p$，可用来回验。

\medskip
\textbf{三问防错：}点在曲线上吗？开口朝哪里？方程系数对应 $2p$ 还是 $4a$？
\end{trapbox}
